\begin{figure}[htbp]
\centering
\begin{tikzpicture}[
  scale=0.6, transform shape, font=\small,
  %%%%%%%%% struct -> clase UML
  structm/.style={draw, fill=white, rectangle split, rectangle split parts=3,
                  rectangle split part align=center, align=left, inner sep=4pt},
  %%%%%%%%% enum -> enumeración UML
  enumer/.style={draw, fill=black!4, rectangle split, rectangle split parts=2,
                 rectangle split part align=center, align=left, inner sep=4pt},
  %%%%%%%%% trait -> interface UML
  trait/.style={draw, dashed, fill=black!4, rectangle split, rectangle split parts=2,
                rectangle split part align=center, align=left, inner sep=4pt},
  %%%%%%%%% relacións
  depende/.style={dashed, -{Latex[open, length=2.4mm, width=2mm]}},
  realiza/.style={dashed, -{Triangle[open, length=3mm, width=3mm]}},
  etiq/.style={font=\scriptsize, inner sep=1.5pt, fill=white},
  membro/.style={font=\footnotesize\ttfamily}
]

%%%%%%%%% o compoñente
\node[structm, minimum width=64mm] (joy) at (0,0) {
  \nodepart{one}   \texttt{Joypad}
  \nodepart[membro]{two}
        + input: u8\\
        - joyp: u8\\
        - previous\_lines: u8
  \nodepart[membro]{three}
        + button\_down(btn, is\_down)\\
        \textasciitilde{} step(interrupt)\\
        - lines()\\
        + read(address) / write(address, value)
};

%%%%%%%%% os botóns
\node[enumer, minimum width=34mm] (btn) at (-62mm,-46mm) {
  \nodepart{one}  \scriptsize$\ll$enum$\gg$\\[1pt] \texttt{JoypadButton}
  \nodepart[membro]{two}
        Right / Left / Up / Down\\
        A / B / Select / Start
};

%%%%%%%%% direccionamento
\node[trait, minimum width=42mm] (acc) at (58mm,-46mm) {
  \nodepart{one}  \scriptsize$\ll$trait$\gg$\\[1pt] \texttt{Accessible<u16>}
  \nodepart[membro]{two}
        + read(address)\\
        + write(address, value)
};

%%%%%%%%% relacións
\draw[depende] (joy.south west) -- node[etiq, above, pos=0.55]{$\ll$emprega$\gg$} (btn.north east);
\draw[realiza] (joy.south east) -- node[etiq, above, pos=0.55]{implementa} (acc.north west);

\end{tikzpicture}
\caption{Estrutura do mando e da súa relación co espazo de direccións.}
\label{fig:joypad-uml}
\end{figure}
